\section{Bài tập 3}

Cho mạch điện như hình bên dưới. Hãy sắp xếp (vẽ) lại mạch điện để làm rõ cấu trúc nối tiếp và/hoặc song song của mạch. Tiếp theo, vận dụng kiến thức đã học để tìm giá trị điện trở tương đương giữa hai đầu $A$ và $F$, cùng các giá trị điện áp tại các điểm $A, B, C, D$ và $E$. Cuối cùng, thực hiện mô phỏng để kiểm tra kết quả tính toán.

\given{ex3_given}{0.65\textwidth}{Tìm điện trở tương đương toàn mạch và điện áp tại $A, B, C, D, E$}

\subsection{Vẽ lại mạch điện}
\drawfig{ex3_redraw}{Mạch điện sau khi vẽ lại dạng hình thang}

\subsection{Tính toán}
\noindent Chọn nút $F$ làm mốc điện thế (đất, $V_F = 0\text{ V}$):
\begin{enumerate}[leftmargin=2em]
  \item \textbf{Tính điện trở tương đương (rút gọn từ phải sang trái):}
        \begin{align*}
        R_{45} &= R_4 + R_5 = 1\text{ k}\Omega + 2\text{ k}\Omega = 3\text{ k}\Omega \\
        R_{DE} &= R_9 \parallel R_{45} = \frac{6 \cdot 3}{6 + 3} = 2\text{ k}\Omega \\
        R_{CE} &= R_6 \parallel (R_3 + R_{DE}) = \frac{6 \cdot (10 + 2)}{6 + (10 + 2)} = \frac{72}{18} = 4\text{ k}\Omega \\
        R_{BE} &= R_{10} \parallel (R_2 + R_{CE}) = \frac{6 \cdot (2 + 4)}{6 + (2 + 4)} = \frac{36}{12} = 3\text{ k}\Omega \\
        R_{BF} &= R_8 \parallel (R_{BE} + R_7) = \frac{4 \cdot (3 + 9)}{4 + (3 + 9)} = \frac{48}{16} = 3\text{ k}\Omega \\
        R_{AF} &= R_1 + R_{BF} = 2\text{ k}\Omega + 3\text{ k}\Omega = 5\text{ k}\Omega
        \end{align*}
  \item \textbf{Tính điện áp các nút ($I_1 = 18\text{ mA}, V_F = 0\text{ V}$):}
        \begin{align*}
        V_A &= I_1 \cdot R_{AF} = 18\text{ mA} \times 5\text{ k}\Omega = 90\text{ V} \\
        V_B &= I_1 \cdot R_{BF} = 18\text{ mA} \times 3\text{ k}\Omega = 54\text{ V} \\
        I_7 &= \frac{V_B}{R_{BE} + R_7} = \frac{54\text{ V}}{3\text{ k}\Omega + 9\text{ k}\Omega} = 4{,}5\text{ mA} \\
        V_E &= I_7 \cdot R_7 = 4{,}5\text{ mA} \times 9\text{ k}\Omega = 40{,}5\text{ V} \\
        I_2 &= \frac{V_B - V_E}{R_2 + R_{CE}} = \frac{54\text{ V} - 40{,}5\text{ V}}{2\text{ k}\Omega + 4\text{ k}\Omega} = \frac{13{,}5\text{ V}}{6\text{ k}\Omega} = 2{,}25\text{ mA} \\
        V_C &= V_B - I_2 R_2 = 54\text{ V} - 2{,}25\text{ mA} \times 2\text{ k}\Omega = 49{,}5\text{ V} \\
        I_3 &= \frac{V_C - V_E}{R_3 + R_{DE}} = \frac{49{,}5\text{ V} - 40{,}5\text{ V}}{10\text{ k}\Omega + 2\text{ k}\Omega} = \frac{9\text{ V}}{12\text{ k}\Omega} = 0{,}75\text{ mA} \\
        V_D &= V_C - I_3 R_3 = 49{,}5\text{ V} - 0{,}75\text{ mA} \times 10\text{ k}\Omega = 42\text{ V}
        \end{align*}
\end{enumerate}

\begin{dapsobox}
  \begin{itemize}[leftmargin=1.5em]
    \item $R_{AF} = 5\text{ k}\Omega$
    \item $V_A = 90\text{ V}, \quad V_B = 54\text{ V}, \quad V_C = 49{,}5\text{ V}, \quad V_D = 42\text{ V}, \quad V_E = 40{,}5\text{ V}$
  \end{itemize}
\end{dapsobox}

\subsection{Mô phỏng}
\simresult
\simfig{ex3_sim}{Kết quả mô phỏng PSpice của Bài tập 3}
